% =====================================================================
% SECTION VII — DISCUSSION (tightened for 4–5 page target)
% =====================================================================

\section{Discussion}
\label{sec:discussion}

\textbf{Alpha as a structural diagnostic.}
The learned $\alpha$ is not merely a performance knob---it is a
dataset-level diagnostic. On SupplyGraph, $\alpha$ converges to
${\approx}0.92$ within four epochs and narrows to a tight band
(Fig.~\ref{fig:alpha}a), confirming temporal dominance on daily
data with rich sequential structure. On USGS, $\alpha$ oscillates
near $0.5$ throughout training (Fig.~\ref{fig:alpha}b), confirming
that 15 typed inter-commodity structural relations contribute as
much as the sparse annual time series. The 43-point gap between
the two converged values is not a hyperparameter outcome but a
data-driven discovery. A practitioner can read this trajectory
post-training to determine whether more temporal history or more
structural edge data would yield greater performance return.

\textbf{Fixed-fusion failure mode.}
Fixed Fusion ($\alpha{=}0.5$) achieves competitive USGS mean WAPE
but with $5\times$ the variance of the adaptive model
(${\pm}5.17\%$ vs $.{\pm}0.83\%$)---the defining failure mode of
fixed weighting: proximity to the true optimum on average cannot
prevent high seed-to-seed instability because the model cannot
express confidence in its blending assumption.

\textbf{Limitations.}
(1)~USGS has only five annual observations per node; $n{=}5$
prevents Wilcoxon from achieving $p{<}0.0625$.
(2)~The 222-day SupplyGraph window constrains the test partition
to the EEA group's late-period behaviour.
(3)~Layer~5 lead times for SupplyGraph are synthetically
generated~\cite{chopra2007supply}; the risk interface is a
proof-of-concept and is not validated against operational
procurement outcomes.
(4)~All experiments use a single T4 GPU; distributed training
characteristics are not reported.

\smallskip
\noindent\textbf{Disclosure.} No private procurement data were
used. Layer~5 should not be used for consequential procurement
decisions without domain-expert validation.
